\documentclass[11pt,a4paper]{article}

\usepackage[utf8]{inputenc}
\usepackage[T1]{fontenc}
\usepackage{lmodern}
\usepackage[margin=2.2cm]{geometry}
\usepackage{graphicx}
\usepackage{booktabs}
\usepackage{array}
\usepackage{longtable}
\usepackage{tabularx}
\usepackage{enumitem}
\usepackage{xcolor}
\usepackage{hyperref}
\usepackage{fancyhdr}
\usepackage{titlesec}
\usepackage{float}
\usepackage{setspace}

\definecolor{Accent}{RGB}{60,60,60}

\hypersetup{
    colorlinks=true,
    linkcolor=Accent,
    urlcolor=Accent,
    citecolor=Accent
}

\setlength{\parindent}{0pt}
\setlength{\parskip}{0.7em}
\onehalfspacing

\titleformat{\section}
  {\Large\bfseries}
  {\thesection.}
  {0.6em}
  {}

\titleformat{\subsection}
  {\large\bfseries}
  {\thesubsection}
  {0.6em}
  {}

\pagestyle{fancy}
\fancyhf{}
\lhead{Under the Sofa}
\rhead{Game Design Document}
\cfoot{\thepage}
\setlength{\headheight}{13.6pt}

\begin{document}
\hypersetup{pageanchor=false}

% ============================================================
% COVER
% ============================================================

\begin{titlepage}
    \centering

    \vspace*{2cm}

    {\Huge\bfseries UNDER THE SOFA\par}

    \vspace{0.5cm}

    {\Large Game Design Document\par}

    \vspace{1.5cm}

    \rule{\textwidth}{0.6pt}

    \vspace{1.5cm}

    \fbox{
        \parbox[c][7cm][c]{0.85\textwidth}{
            \centering
            \includegraphics[width=\linewidth,height=7cm,keepaspectratio]{screenshot.png}
        }
    }

    \vfill

    \begin{tabular}{ll}
        \textbf{Document Version:} & 1.0 \\[0.3em]
        \textbf{Project:} & Under the Sofa \\[0.3em]
        \textbf{Developer:} & Solo independent project \\[0.3em]
        \textbf{Target Platform:} & Windows PC \\[0.3em]
        \textbf{Engine:} & Unity 6.5 \\[0.3em]
        \textbf{Render Pipeline:} & Universal Render Pipeline \\[0.3em]
        \textbf{Date:} & September 2026 \\[0.3em]
        \textbf{Author:} & Matej Spaji\'c \\[0.3em]
        \textbf{Repository:} & \href{https://github.com/Matej293/UnderTheSofa}{github.com/Matej293/UnderTheSofa} \\
    \end{tabular}

    \vfill

    {\small This document describes the current playable game.}

\end{titlepage}

\hypersetup{pageanchor=true}
\pagenumbering{roman}
\tableofcontents
\clearpage
\pagenumbering{arabic}

% ============================================================
% CURRENT GAME DESIGN
% ============================================================

\section{Game Summary}\label{game-summary}

\emph{Under the Sofa} is a single-player driving and exploration game about a small RC car inside an oversized house. The player drives through a connected indoor course, climbs household objects, searches for routes, and collects rings.

The game uses arcade Rigidbody physics. Handling is light, fast, and a little unstable. The car can drift, hop, use a rechargeable speed boost, and rotate in the air. Missed jumps are cheap because the player can rewind or move back to a safe position.

The player advances by collecting rings and learning how the car moves.

\subsection{Main Player Thought}\label{main-player-thought}

\begin{quote}
How do I get up there?
\end{quote}

Furniture and props form the driving course. A book can become a ramp, a table can become a platform, and a gap between household objects can become a jump. Rooms favor readable movement routes over strict realism.

\subsection{Creative Reference}\label{creative-reference}

\emph{Re-Volt} (1999) is the main reference for vehicle feel and atmosphere. \emph{Under the Sofa} borrows the sense of driving a fast toy car through large everyday spaces. The game uses that movement for exploration and collecting rather than conventional racing.

The audio and interface follow a late-1990s electronic game style. Low-resolution presentation, simple geometry, small textures, direct menus, and trance music support the same identity.

\section{Current Playable Scope}\label{current-playable-scope}

The Windows build contains two scenes:

\begin{itemize}
\item
  \texttt{MainMenu}, a live title scene set inside the house
\item
  \texttt{Vehicle}, the complete playable course
\end{itemize}

The gameplay scene is one connected indoor environment. Its named sections include the living room, bathroom, bedroom, and an extra-room section. The bathroom and bedroom are tied to the two ring gates.

The full run uses one car and 17 rings. A first door opens at 7 rings. A second door opens at 12 rings. Collecting all 17 rings completes the game.

\section{Player Experience}\label{player-experience}

Driving remains active even when the player is not following the ring route. The car reacts to ramps, furniture edges, uneven surfaces, and collisions. The player can look for a cleaner approach, attempt a larger jump, or use air control to repair a poor landing.

The ring path gives the course a clear goal. It also draws the player into different parts of the house and toward elevated objects. Progress comes from the same movement used during free exploration.

Failure costs little. \texttt{R} rewinds the car through its recent movement, \texttt{U} returns it to an earlier safe pose, and a fall below the course triggers an automatic recovery.

\section{Gameplay Loop}\label{gameplay-loop}

The player starts in the house and begins driving. Rings provide the immediate route through the environment. Furniture and loose props create ramps, ledges, narrow paths, and places to build speed.

A normal sequence looks like this:

\begin{quote}
Find a ring route, line up the car, build speed, jump, correct the car in the air, land, and continue.
\end{quote}

The first seven rings open the bathroom route. Five more rings open the bedroom route. The remaining five lead to game completion. The counter is cumulative across the whole scene.

After completion, the player can continue exploring or return to the main menu.

\section{Vehicle}\label{vehicle}

\subsection{Player Car}\label{player-car}

The project has one playable car. Its visual model is an Acura Integra Type R adapted as the RC vehicle.

The car keeps the same abilities and tuning for the whole game. Rings open doors through the connected course.

\subsection{Physics Setup}\label{physics-setup}

The vehicle uses a standard Unity Rigidbody with custom arcade forces. The current player prefab uses these main Rigidbody values:

\begin{longtable}[]{@{}lr@{}}
\toprule\noalign{}
\textbf{Setting} & \textbf{Current Value} \\
\midrule\noalign{}
\endhead
\bottomrule\noalign{}
\endlastfoot
Mass & 1.6 \\
Linear damping & 0.15 \\
Angular damping & 0.35 \\
Gravity & On \\
Interpolation & On \\
Collision detection & Continuous Dynamic \\
\end{longtable}

The controller applies drive force near the rear axle. Separate front and rear lateral grip keep the car readable while still allowing slides. Ground checks use a sphere cast and accept slopes up to 55 degrees.

The center of mass is lowered and an upright force steadies the car while grounded. A low-edge assist gives a small upward push when the front of the car catches a short obstacle. This prevents minor collider edges from stopping normal movement.

\subsection{Driving}\label{driving}

\texttt{W} accelerates and \texttt{S} brakes or reverses. Steering works in both directions and changes with speed. Normal drive force uses a forward threshold of 22 Unity units per second and a reverse threshold of 9.

The wheels turn and spin visually. Their motion follows steering input and forward speed.

\subsection{Drift}\label{drift}

Holding either Shift key while steering starts a drift once the car has enough speed. Drift lowers rear grip and applies handbrake deceleration. Normal drive and boost force are cut while the handbrake is held.

Drifting helps rotate the car through tight approaches. The slide still comes from Rigidbody motion, so speed and entry angle affect the result.

\subsection{Hop}\label{hop}

Space applies a short upward impulse while the car is grounded. The hop clears small edges and helps with quick corrections. Ground detection enables the hop after the car lands.

Ramps and household objects provide the main height for large jumps.

\subsection{Air Control}\label{air-control}

The player can rotate the car while airborne:

\begin{itemize}
\item
  \texttt{W} and \texttt{S} control pitch
\item
  \texttt{A} and \texttt{D} control yaw
\item
  \texttt{Q} and \texttt{E} control roll
\end{itemize}

These controls allow flips, rolls, and landing corrections.

\subsection{Boost}\label{boost}

Holding the Ctrl key uses the speed boost. The meter starts full and holds four seconds of energy. Energy drains at one unit per second and recharges at 0.6 units per second whenever Ctrl is released.

Boost adds forward acceleration and produces orange particles behind the car. The HUD meter also changes while boost is active. The meter starts full at the beginning of each run.

\subsection{Vehicle Effects}\label{vehicle-effects}

The car emits light dust while moving on the ground. Boost has a separate orange particle effect. These effects scale with vehicle speed or boost state.

Crashes affect the car\textquotesingle s motion and can be corrected with the recovery controls.

\section{Recovery}\label{recovery}

Recovery is part of normal play. The current game has three related systems.

\subsection{Rewind}\label{rewind}

Pressing \texttt{R} rewinds up to three seconds of recent position and rotation history. Playback lasts 1.5 seconds. The controller locks during the rewind, and vehicle collisions are disabled until playback ends.

Rewind follows the recorded path. It is the quickest way to retry a missed approach without returning to the start of the room.

\subsection{Unstuck}\label{unstuck}

Pressing \texttt{U} searches the recent safe-pose history for a point at least 1.5 units away. The chosen pose must be on sufficiently level ground. The car is placed above that point with its motion cleared.

If no suitable recent pose exists, the system falls back to the initial safe pose.

\subsection{Fall Recovery}\label{fall-recovery}

The car records safe grounded poses during play. Falling below a height of -8 resets it to the latest safe pose and clears linear and angular velocity.

\section{Controls}\label{controls}

The game targets keyboard play on Windows.

\subsection{Driving Controls}\label{driving-controls}

\begin{longtable}[]{@{}ll@{}}
\toprule\noalign{}
\textbf{Input} & \textbf{Action} \\
\midrule\noalign{}
\endhead
\bottomrule\noalign{}
\endlastfoot
W & Accelerate, or pitch forward in the air \\
S & Brake and reverse, or pitch backward in the air \\
A / D & Steer on the ground and control yaw in the air \\
Space & Hop while grounded \\
Left Shift & Drift and handbrake \\
Left Ctrl & Boost \\
Q / E & Air roll \\
R & Rewind recent movement \\
U & Move to an earlier safe pose \\
Esc & Pause \\
\end{longtable}

\subsection{Menu Controls}\label{menu-controls}

\begin{longtable}[]{@{}ll@{}}
\toprule\noalign{}
\textbf{Input} & \textbf{Action} \\
\midrule\noalign{}
\endhead
\bottomrule\noalign{}
\endlastfoot
Up / Down Arrow & Move through menu entries \\
Left / Right Arrow & Adjust the selected option \\
Enter & Confirm \\
Esc & Go back, resume, or cancel a quit prompt \\
\end{longtable}

\section{Camera}\label{camera}

The game uses a fixed third-person chase camera. It follows the car automatically and keeps world up as its horizon reference, even while the car flips or rolls.

At low speed, the camera sits 4.5 units behind the car. It moves toward a 7.5-unit distance as speed approaches 16 units per second. Boost adds one unit of distance, and being airborne adds 1.25 units.

The camera smooths both position and direction. It follows the car\textquotesingle s flattened forward heading instead of copying the car\textquotesingle s full rotation. This keeps the horizon stable and prevents airborne rolls from turning the whole view.

A sphere cast checks the route between the target and the intended camera position. When furniture or a wall blocks the view, the camera moves closer. Its minimum distance is 1.25 units.

\section{Rings and Progression}\label{rings-and-progression}

Rings are the only collectible and the only progression counter. Each ring is worth one point. A ring hides its visual and trigger when collected, plays its pickup feedback, then deactivates.

The gameplay scene contains 17 ring instances. Ring progress is cumulative:

\begin{longtable}[]{@{}rl@{}}
\toprule\noalign{}
\textbf{Ring Count} & \textbf{Result} \\
\midrule\noalign{}
\endhead
\bottomrule\noalign{}
\endlastfoot
7 & Bathroom door opens and the HUD reports the unlock \\
12 & Bedroom door opens and the HUD reports the unlock \\
17 & The game-complete message appears \\
\end{longtable}

Both doors rotate open over one second and play a door sound. The HUD notifications tell the player when each threshold has been reached.

Starting a new game reloads the gameplay scene with all 17 rings ready to collect.

\section{World and Level Design}\label{world-and-level-design}

\subsection{House Course}\label{house-course}

The current game takes place inside one furnished house scene. The environment uses a prebuilt house model with added furniture, office props, bathroom props, toys, a dresser, food items, and other small objects. Custom mesh colliders and improved hitbox objects make the imported geometry driveable.

The living room forms the start of the route. Cumulative ring gates then open access associated with the bathroom and bedroom. All progression happens without changing scenes.

The course uses household scale to turn ordinary objects into driving geometry. Tables, shelves, books, paper, furniture, and clutter provide changes in height and places to jump. The player can keep exploring after collecting the final ring.

\subsection{Route Rules}\label{route-rules}

Every route uses the current car\textquotesingle s fixed movement values.

The hop handles low obstacles. Larger height changes come from the shape and placement of the environment. Ring placement points toward usable paths and progression doors.

Most environment props stay fixed so their collision shapes provide consistent driving routes.

\section{HUD}\label{hud}

The gameplay HUD uses a 640 by 360 reference layout.

\begin{itemize}
\item
  A boost bar sits at the lower left.
\item
  The current ring count sits at the upper right.
\item
  Temporary messages appear near the top center when a door opens or the game ends.
\end{itemize}

The ring panel pulses when the count changes. The boost bar drains and refills smoothly, with small orange particles while boost is active.

\section{Front End}\label{front-end}

\subsection{Main Menu}\label{main-menu}

The game opens in the \texttt{MainMenu} scene. A live house view and the player car sit behind the interface. The menu music fades in, and choosing New Game triggers a short camera move and fade into the \texttt{Vehicle} scene.

The main menu contains:

\begin{itemize}
\item
  New Game
\item
  Options
\item
  Controls
\item
  Credits
\item
  Quit
\end{itemize}

\subsection{Pause Menu}\label{pause-menu}

Pressing Esc during play pauses time, vehicle input, and gameplay music. Losing application focus also opens the pause menu.

The pause menu contains Resume, Options, Controls, Main Menu, and Quit. Returning to play restores time, input, and music.

\subsection{Options}\label{options}

The options page contains four settings:

\begin{itemize}
\item
  Master Volume
\item
  Music Volume
\item
  Fullscreen On or Off
\item
  Resolution
\end{itemize}

Volume changes in ten-percent steps. The resolution list comes from display modes reported by Unity. These four settings persist through later sessions with \texttt{PlayerPrefs}.

\subsection{Controls and Credits}\label{controls-and-credits}

The Controls page lists the main movement and recovery keys. The Credits page names Matej Spaji\'c as the creator.

\subsection{Completion Screen}\label{completion-screen}

Collecting the seventeenth ring shows a completion notification, then pauses play on the Game Complete page. The player can choose Continue Exploring or Main Menu.

\section{Visual Direction}\label{visual-direction}

The game uses low-poly imported models and a late-1990s presentation. The car, furnished house, props, fonts, HUD, and color choices favor clear shapes and strong blocks of color.

Gameplay renders through a 640 by 360 texture with point filtering and no anti-aliasing. A presentation layer scales this image to the display and keeps a 16:9 game area with letterboxing when needed.

The current post-processing profiles include color adjustment, white balance, tonemapping, vignette, motion blur, and film grain.

The Silkscreen font and a cropped retro UI texture define the menu and HUD style. Text stays large enough to read at the internal resolution.

\section{Audio}\label{audio}

\subsection{Music}\label{music}

Music follows the late-1990s trance and acid electronic direction. The current build has one looping menu track and a three-track gameplay playlist.

Starting a new game fades out the menu track. Gameplay tracks play in sequence, fading near the end of one track before the next begins. Pausing stops the current gameplay track and resuming continues it.

\subsection{Sound Effects}\label{sound-effects}

The current sound set covers menu navigation, confirmation, cancellation, errors, notifications, ring collection, and both progression doors.

\section{Technical Setup}\label{technical-setup}

\begin{longtable}[]{@{}ll@{}}
\toprule\noalign{}
\textbf{Component} & \textbf{Current Setup} \\
\midrule\noalign{}
\endhead
\bottomrule\noalign{}
\endlastfoot
Unity & 6000.5.10f1 \\
Render Pipeline & URP 17.5.0 \\
Input System & 1.20.0 \\
Unity Pipeline & 0.5.0-exp.1 \\
Unity CLI & 1.0.0-beta.6 \\
UI & Unity UI 2.5.0 \\
Target & Windows PC \\
Gameplay physics & Rigidbody with custom arcade forces \\
Build scenes & MainMenu and Vehicle \\
Internal game view & 640 by 360, point filtered \\
\end{longtable}

Asset sources and attribution are tracked separately in \texttt{Docs/asset-references.xlsx}. That spreadsheet remains the source for asset reference details.

\section{Design Rules for the Current Game}\label{design-rules-for-the-current-game}

\begin{itemize}
\item
  Keep the single car responsive and readable.
\item
  Build routes around the car\textquotesingle s fixed movement values.
\item
  Use ramps and household geometry for large jumps.
\item
  Keep the hop small and grounded.
\item
  Make failed jumps quick to retry through rewind and unstuck.
\item
  Use rings to lead the player through the connected indoor course.
\item
  Keep all ring gates reachable with the base car.
\item
  Preserve the automatic chase camera and stable horizon.
\item
  Keep the interface small and keyboard-friendly.
\item
  Treat the 17-ring run as the complete progression structure.
\end{itemize}

\clearpage
\renewcommand{\refname}{Assets}
\addcontentsline{toc}{section}{Assets}
\nocite{*}
\bibliographystyle{unsrt}
\bibliography{bibliography}

\end{document}
